\chapter{LITERATURE REVIEW}
\label{ch:lit}

\section{TWO WHEELED INVERTED PENDULUM}

Research on two wheeled inverted pendulums began in the 1990s. Ha and Yuta at
the University of Tsukuba created the Yamabico Kurara and used a linear
quadratic regulator design for trajectory tracking control \cite{ha1994trajectory}.
Shiroma et al.\ built a wheeled inverted pendulum and analyzed its stability in
the presence of applied forces, using an observer to estimate the unknown force
\cite{shiroma1996cooperative}. Ozaki et al.\ implemented decoupling control by
adding an additional state variable, creating a square system, and proposed
algorithms to cope with the singularity introduced by the feedback
\cite{ozaki2001position}. Grasser et al.\ at the Industrial Electronics
Laboratory of the Swiss Federal Institute of Technology created JOE, a mobile
inverted pendulum, decoupling the robot dynamics and applying pole placement
together with a filter to suppress the backlash effects of the DC motors
\cite{grasser2002joe}. Ooi at the University of Western Australia built a
balancing robot to investigate Kalman filtering for sensor fusion, implementing
both pole placement and linear quadratic regulator controllers
\cite{ooi2003balancing}. Various forms of PID control have been investigated by
a number of researchers
\cite{solis2011waseda,sun2010researching,guerrero2013realtime,miranda2009application}.
Kim and Kwon developed a feedforward controller based on the state dependent
Riccati equation \cite{kim2013nonlinear}.

\rev{Two observations about this body of work bear on the present development.
First, every one of these designs depends on the physical parameters being known
accurately, since the gains are computed from them; the sensitivity of the
closed loop to parameter error is the practical limitation rather than the
linearization itself. Second, Grasser et al.\ treat backlash as something to be
filtered out of the measurement rather than compensated in the control, which is
a reasonable engineering choice but leaves the nonlinearity in the loop. Both
points recur in Chapter~\ref{ch:platform}.}

The work above is based on linearization of the equations of motion. Nonlinear
control of two wheeled inverted pendulums is a current topic of research. Kim et
al.\ investigated the exact dynamics of a two wheeled inverted pendulum robot,
examining stability on an inclined plane and in turning motion
\cite{kim2005dynamic}. Pathak et al.\ used the full nonlinear equations obtained
by the Euler-Lagrange method and applied partial feedback linearization to
control velocity and position \cite{pathak2005twip}. \rev{This last is worth
singling out: the partial feedback linearization decomposition makes explicit
that the actuated and unactuated subsystems must be treated differently, and
that the tilt state rather than the motor voltage is the effective control
authority available to the position loop. That observation is the structural
basis for the virtual control used in Chapter~\ref{ch:control}.} Askari et al.\
implemented a form of model predictive control and studied performance under
input disturbances \cite{askari2009model}. Shibayama et al.\ implemented an
observer-based robust controller to counteract unknown disturbances, permitting
the use of a simplified system model \cite{shibayama2010robust}. Ja and Lee
implemented sliding mode control, improving on linear controllers when far from
equilibrium \cite{ja2012position}. Tsai and Ju implemented backstepping sliding
mode control for trajectory tracking and stabilization by first decoupling the
kinematics and dynamics \cite{tsai2010trajectory}. Do and Seet combined partial
feedback linearization with $p$-times differentiable saturation and backstepping;
the resulting controller has a large domain of attraction and is simple to tune
\cite{do2010motion}. Rudra and Barai developed a robust adaptive backstepping
technique that estimates the system parameters and removes the need for prior
knowledge of them \cite{rudra2012robust}. Durdevic and Yang proposed a hybrid
switching controller to overcome the backlash nonlinearity of DC motors and
demonstrated its effectiveness experimentally \cite{durdevic2013hybrid}. Yue,
Wei, and Li developed an adaptive sliding mode technique based on zero dynamics
theory, allowing parameter uncertainties to be estimated and providing
robustness in the presence of nonlinearities \cite{yue2014adaptive}.

The general backstepping machinery underlying several of these designs is due to
Krsti\'c et al.\ \cite{krstic1995nonlinear}, and the standard stability tools
used throughout are collected in Khalil \cite{khalil2002nonlinear} and in Slotine
and Li \cite{slotine1991applied}.

Neural networks have been used by a number of researchers to control the two
wheeled inverted pendulum. Noh, Lee, and Jung used radial basis function
networks to control the unstable system, deriving an online training law from
the back-propagation algorithm \cite{noh2008position,jung2008control}. Tsai,
Huang, and Lin proposed an adaptive technique using radial basis function
networks, employing backstepping and Lyapunov analysis to synthesize stable
adaptive control laws \cite{tsai2010adaptive}. Li and Yang formulated an output
feedback adaptive neural network controller with a linear dynamic compensator
and showed that it outperforms a model based controller when parametric
uncertainties are present \cite{li2012neural}. Li, Yang, and Fan devote a
chapter of their book on advanced nonlinear control of wheeled inverted
pendulums to neural network control using radial basis functions
\cite{li2013advanced}. Optimal control remains a smaller area of research for
these systems, though Gomez et al.\ developed a strategy based on Control
Adjoining Cell Mapping Reinforcement Learning and tested it on a robot built
from LEGO NXT components \cite{gomez2012optimal}.

Applications of the wheeled inverted pendulum configuration extend beyond the
laboratory; Neugebauer et al.\ describe mobile systems of this type used for
machining large work pieces \cite{neugebauer2012mobile}.

\rev{The common difficulty across all of the nonlinear and adaptive work above,
for the TWIP specifically, is underactuation. A single input drives both
acceleration equations, so an uncertainty appearing in the tilt acceleration
cannot be cancelled by the same signal that cancels an uncertainty in the
translational acceleration. Two of the works cited treat the actuator
nonlinearity directly rather than absorbing it into a lumped uncertainty
\cite{grasser2002joe,durdevic2013hybrid}, and the broader adaptive treatment of
actuator nonlinearities is due to Tao and Kokotovi\'c \cite{tao1993backlash}.
This matters here because backlash is not smooth and therefore falls outside the
function class covered by the universal approximation property. A network can
approximate such a nonlinearity only in an averaged sense, and the residual
appears as a floor on the achievable tracking error rather than as something
further adaptation can remove; see Remark~\ref{rem:backlash}.}

\rev{\subsection{Command Filtering}

Classical backstepping requires analytic differentiation of each virtual control
law, which produces an expanding chain of terms and, when the virtual control
contains an online function estimate, requires differentiating the weight update
law itself. Command filtering, introduced by Farrell et al.\
\cite{farrell2009command} and developed as dynamic surface control by Swaroop et
al.\ \cite{swaroop2000dsc}, replaces analytic differentiation with a
bandwidth-limited filter and compensates for the resulting filter error within
the stability argument. This device is used in Chapter~\ref{ch:control}, where
it serves a second purpose beyond avoiding the explosion of terms: it breaks the
algebraic dependence that would otherwise exist between the two neural networks'
approximation targets, as detailed in Remark~\ref{rem:order}.}

\section{MODIFIED STATE OBSERVER}

Artificial neural networks have emerged as a capable tool in adaptive control.
Modeled after the behavior of neurons in biological systems, they have proven
effective as function approximators \cite{hornik1989multilayer,sadegh1993nonlinear}.
This approximation capability has been adapted into many methods of adaptive
control that use networks to estimate nonlinear system uncertainty
\cite{yang2013electrohydraulic,jagannathan2006neural}. Rajagopal et al.\
developed the Modified State Observer (MSO) as a method of estimating nonlinear
uncertainty by combining a standard Luenberger observer structure with a neural
network \cite{rajagopal2009mso}. The MSO allows the designer to use large
adaptive gains without the high frequency oscillation in the control signal that
could otherwise excite unmodeled dynamics.

Although originally applied to control, the MSO has also been used purely for
state estimation. Harl et al.\ applied it to an orbit uncertainty estimation
problem and extended the method to a reduced order formulation for cases where
the full state cannot be measured \cite{harl2011orbit}. Darling et al.\ used the
MSO and its reduced order form in an atmospheric reentry uncertainty estimation
problem, estimating the uncertain aerodynamic acceleration of falling debris
\cite{darling2012reentry}, and subsequently developed the Sigma Point MSO, which
applies sigma point filtering in the manner of the unscented Kalman filter and
admits nonlinear measurements with a variable observer gain
\cite{darling2013sigma}. Yang et al.\ formulated a model reference adaptive
control design based on the MSO to control a nonlinear electrohydraulic system
with parameter uncertainty from a simplified linear model
\cite{yang2013electrohydraulic}. Pappu et al.\ formulated an MSO-based adaptive
control law and applied it to the Black Kite micro air vehicle
\cite{pappu2014blackkite}.

All of this previous work on the MSO is in continuous time. In practice,
implementation on a digital system requires integrating the dynamics over time
at every step. Discrete time implementations have the advantage of reduced
computation in addition to being designed for digital hardware directly.
Discrete time neural networks have themselves been a topic of study in adaptive
control \cite{jagannathan2006neural}. To a lesser extent they have been used
specifically as observers. Salgado and Chairez used recurrent networks with
offline training based on a least mean square method to construct a discrete
time neural observer \cite{salgado2013nonlinear}. Alanis et al.\ proposed a
reduced order discrete time neural observer trained offline by an extended
Kalman filter using high order recurrent networks \cite{alanis2010reduced}. What
these works lack is an online training law with guaranteed boundedness. Lewis et
al.\ developed a multilayer discrete time neural network observer using online
training \cite{lewis1999neural}; the present work differs in the weight update
law and in the observer structure, and requires substantially fewer tunable
parameters, which simplifies both the design and the tuning effort.

\rev{\subsection{Method Notes}

Three points of technique motivate the proof given in Chapter~\ref{ch:dmso}.

First, the weight error recursion in discrete time adaptive schemes has the form
$\Wt_{k+1}=(I-\Gamma\phi_k\phi_k^T)\Wt_k+(\cdot)$, in which
$I-\Gamma\phi\phi^T$ is a symmetric rank-one perturbation of the identity. Its
spectrum is $\{1,\,1-\Gamma\norm{\phi}^2\}$, so its induced two-norm is
$\max\{1,\abs{1-\Gamma\norm{\phi}^2}\}$ and not $\abs{1-\Gamma\norm{\phi}^2}$
\cite{horn2013matrix}. Bounding this quantity before the Lyapunov difference is
formed discards the structure that the adaptation law was constructed to
produce.

Second, the condition $0<\Gamma\norm{\phi_k}^2<2$ familiar from the persistency
of excitation literature \cite{sadegh1993nonlinear,lewis1999neural} is recovered
directly from the exact first difference, and is not an independent assumption.

Third, weight drift in the absence of persistent excitation is a long-standing
concern in adaptive control, and $\sigma$-modification was introduced by Ioannou
and Kokotovi\'c \cite{ioannou1983adaptive,ioannou1984sigma} precisely to address
it. The device is used twice in this work: on the controller networks in
Chapter~\ref{ch:control}, as in the previous revision, and additionally on the
observer network in Section~\ref{sec:sigmod}. Applying it in both places makes
the two stability arguments structurally consistent and removes reliance on an
excitation condition that regulation about an equilibrium does not supply.}

\section{SUPPORTING MATERIAL}

The matrix algebra used in Chapter~\ref{ch:dmso}, in particular the properties of
semi-orthogonal matrices and of the Frobenius norm, follows Abadir and Magnus
\cite{abadir2005matrix} and Horn and Johnson \cite{horn2013matrix}. Young's
inequality is due to Young \cite{young1912classes}, and the triangle inequality
and related numerical results follow Kress \cite{kress1998numerical}. The
rotation sequences used to obtain the tilt angle from the accelerometer in
Section~\ref{sec:measurements} follow Kuipers \cite{kuipers1999quaternions}, and
the complementary filter used to fuse the accelerometer and gyroscope follows
Higgins \cite{higgins1975comparison}. The rigid body dynamics used in
Chapter~\ref{ch:model} follow Hibbeler \cite{hibbeler2010dynamics}, and the DC
motor and mechatronic subsystem models follow de Silva
\cite{desilva2004mechatronics}. The virtual control construction extended in
Chapter~\ref{ch:control} is due to Huang \cite{huang2002lyapunov}. The Kalman
filter baseline follows Kalman \cite{kalman1960} and Simon
\cite{simon2006optimal}, and the inertial sensor error models used in
Section~\ref{sec:noise} follow Gebre-Egziabher \cite{gebreegziabher2004design}
and Flenniken \cite{flenniken2005modeling}.
